%% Appendix: the semantics of ISWIMμ and x86μ in the other four styles, and
%% the correctness theorem of each style. Sources:
%% agda/Mu/Theorems/{ScottDomains,AbstractMachines,ClockedInterpreters,ReductionSemantics}.agda (the
%% statements); agda/Mu/Semantics/ClockedInterpreters/{Values,ISWIM,X86}.agda,
%% agda/Framework/Clocked.agda; agda/Mu/Semantics/AbstractMachines/{Values,ISWIM,X86}.agda,
%% agda/Framework/Smallstep.agda; agda/Mu/Semantics/ReductionSemantics/ISWIM.agda,
%% agda/Framework/Reduction.agda (red is shown with step-in unfolded);
%% agda/Framework/Scott/{Domain,Ops,Semantics}.agda,
%% agda/Semantics/Scott/Prim.agda, agda/Mu/Semantics/Scott/{ISWIM,X86}.agda.

\raggedbottom
\allowdisplaybreaks
\section{The Other Four Semantics and Their Theorems}
\label{app:other}
\agda{Mu/Theorems/ClockedInterpreters.agda, AbstractMachines.agda, ReductionSemantics.agda, ScottDomains.agda}

Each of the five developments of Section~\ref{sec:evaluation} proves
the whole compiler correct with respect to its own semantics of the
languages. For the graph model, the semantics are those of
Appendix~\ref{app:languages} and the theorem is
Theorem~\ref{thm:compiler}. This appendix gives, for each of the other
four styles, the semantics of the source language \iswim\ and the
target language x86$\mu$, and the correctness theorem for the whole
compiler, so that each theorem can be read on its own. The semantics
of the intermediate languages follow the same pattern and are in the
artifact.

\subsection{Values and Observations}
\label{app:observations}
\agda{Semantics/Observation.agda, Framework/Smallstep.agda, Framework/Clocked.agda}

The three operational styles share their values and their
observations. A value of type $A$ is one of
\[
  v, w ::= c \mid \kw{code}\,f \mid \kw{pack}(v, w) \mid \kw{tup}(v_1, \ldots, v_n)
    \mid \kw{inl}\,v \mid \kw{inr}\,v \mid \kw{fold}\,v
\]
where a closure $\kw{pack}(v, w)$ pairs a code value $v$ with the tuple
$w$ of its environment. In \iswim, the code of a closure is the body $N$
of a $\lambda$-abstraction and its environment is the environment
$\rho$ of the $\lambda$-abstraction, a list of values; we write
$\langle \lambda x.\,N, \rho\rangle$ for $\kw{pack}(\kw{code}\,N,
\kw{tup}(\rho))$. In x86$\mu$, the code of a closure is the name $f$ of a
definition. A run is observed through its final value, forgetting the
code and environment inside closures. An observation is one of
\[
  o ::= c \mid \mathit{fun} \mid \mathit{tuple}(o_1, \ldots, o_n)
    \mid \mathit{left}(o) \mid \mathit{right}(o) \mid \mathit{error}
\]
and the observation of a value is
\[
\begin{array}{@{}r@{\;}c@{\;}l@{\qquad}r@{\;}c@{\;}l@{}}
  \mathit{obs}(c) & = & c &
  \mathit{obs}(\kw{pack}(v, w)) & = & \mathit{fun}\\
  \mathit{obs}(\kw{tup}(v_1, \ldots, v_n)) & = & \mathit{tuple}(\mathit{obs}(v_1), \ldots, \mathit{obs}(v_n)) &
  \mathit{obs}(\kw{fold}\,v) & = & \mathit{obs}(v)\\
  \mathit{obs}(\kw{inl}\,v) & = & \mathit{left}(\mathit{obs}(v)) &
  \mathit{obs}(\kw{inr}\,v) & = & \mathit{right}(\mathit{obs}(v))
\end{array}
\]
and $\mathit{obs}(\kw{code}\,f) = \mathit{fun}$. A run stops
with the observation $\mathit{error}$ when it gets stuck, which no well-typed program
reaches. The primitives compute $\delta_{\kw{sub}}(n, m) = n - m$ and
$\delta_{\kw{eq}}(n, m) = \kw{inl}\,\kw{tup}()$ if $n = m$ and
$\kw{inr}\,\kw{tup}()$ otherwise.

\paragraph{Machines}

A machine is a set of states, a start state, and a step function that
maps each state $\varsigma$ either to a next state, written $\varsigma \longrightarrow \varsigma'$, or to a final observation $o$, written $\varsigma \longrightarrow \mathsf{stop}(o)$. A machine \emph{terminates with} $o$ when a sequence of steps leads from the start state to a state $\varsigma$ with $\varsigma \longrightarrow \mathsf{stop}(o)$, and it \emph{diverges} when
for every $n$ there is a sequence of $n$ steps from the start state. We use the same notation for each machine; the CEK
machines of \iswim\ and x86$\mu$ (Appendix~\ref{app:machines}) and
the reduction semantics of \iswim\ (Appendix~\ref{app:reduction})
each define their own $\longrightarrow$ and $\longrightarrow
\mathsf{stop}$.

\paragraph{Clocked interpreters}

A clocked interpreter maps a clock $k \in \mathbb{N}$ to a
\emph{result}, ranged over by $e$, which times out, reports an error,
or returns a value $v$:
\[
  e ::= \mathsf{timeout} \mid \mathsf{error} \mid \mathsf{ret}\,v
\]
We write $v \leftarrow e;\; e'$ for the result of $e'$
when $e$ is $\mathsf{ret}\,v$, and for $e$ itself when $e$ is
$\mathsf{timeout}$ or $\mathsf{error}$. An interpreter \emph{terminates
with} $o$ when for some clock $k$ its result is $\mathsf{ret}\,v$ with
$\mathit{obs}(v) = o$, or $\mathsf{error}$ with $o = \mathit{error}$,
and it \emph{diverges} when its result is $\mathsf{timeout}$ for every
clock.

\paragraph{Equivalence}

Two machines, or two interpreters, are \emph{equivalent} when, for
every observation $o$, one terminates with $o$ if and only if the other
does, and one diverges if and only if the other does. The three
operational theorems state that a program and its compilation are
equivalent.

\subsection{Clocked Interpreters}
\label{app:interp}
\agda{Mu/Semantics/ClockedInterpreters/ISWIM.agda, Mu/Semantics/ClockedInterpreters/X86.agda}

\paragraph{Interpreter for \iswim}

The interpreter $\mathit{eval}_k(\rho, M)$ evaluates a term $M$ in an
environment $\rho$ of values at clock $k$. The clock ticks at each call
of the interpreted program.

\defhead{The interpreter of \iswim}{\mathit{eval}_k(\rho, M) = e \qquad \mathit{call}_k(f, w) = e}
\begin{align*}
  \mathit{eval}_k(\rho, x) &= \mathsf{ret}\,\rho(x) \\
  \mathit{eval}_k(\rho, \lambda x.\,N) &= \mathsf{ret}\,\langle \lambda x.\,N, \rho\rangle \\
  \mathit{eval}_k(\rho, L\;M) &= f \leftarrow \mathit{eval}_k(\rho, L);\; w \leftarrow \mathit{eval}_k(\rho, M);\; \mathit{call}_k(f, w) \\
  \mathit{eval}_k(\rho, c) &= \mathsf{ret}\,c \\
  \mathit{eval}_k(\rho, \kw{tuple}\,(M_1, \ldots, M_n)) &= v_1 \leftarrow \mathit{eval}_k(\rho, M_1);\; \ldots;\; v_n \leftarrow \mathit{eval}_k(\rho, M_n);\;
    \mathsf{ret}\,\kw{tup}(v_1, \ldots, v_n) \\
  \mathit{eval}_k(\rho, \kw{get}_i\,M) &= \kw{tup}(v_1, \ldots, v_n) \leftarrow \mathit{eval}_k(\rho, M);\; \mathsf{ret}\,v_i \\
  \mathit{eval}_k(\rho, \kw{inl}\,M) &= v \leftarrow \mathit{eval}_k(\rho, M);\; \mathsf{ret}\,(\kw{inl}\,v)
  \qquad\text{(likewise $\kw{inr}$ and $\kw{fold}$)} \\
  \mathit{eval}_k(\rho, \kw{case}\;L\;(x.\,M)\;(y.\,N)) &= u \leftarrow \mathit{eval}_k(\rho, L);\;
    \begin{cases}
      \mathit{eval}_k(\rho[x{:=}v], M) & \text{if } u = \kw{inl}\,v\\
      \mathit{eval}_k(\rho[y{:=}v], N) & \text{if } u = \kw{inr}\,v
    \end{cases} \\
  \mathit{eval}_k(\rho, \kw{prim}_p\;L\;M) &= n \leftarrow \mathit{eval}_k(\rho, L);\; m \leftarrow \mathit{eval}_k(\rho, M);\; \mathsf{ret}\,\delta_p(n, m) \\
  \mathit{eval}_k(\rho, \kw{let}\;x = M\;\kw{in}\;N) &= v \leftarrow \mathit{eval}_k(\rho, M);\; \mathit{eval}_k(\rho[x{:=}v], N) \\
  \mathit{eval}_k(\rho, \kw{unfold}\;M) &= \kw{fold}\,v \leftarrow \mathit{eval}_k(\rho, M);\; \mathsf{ret}\,v \\[4pt]
  \mathit{call}_0(f, w) &= \mathsf{timeout} \\
  \mathit{call}_{k+1}(\langle \lambda x.\,N, \rho\rangle, w) &= \mathit{eval}_k(\rho[x{:=}w], N)
\end{align*}
A pattern on the left of $\leftarrow$ always matches, because the
interpreter is defined on well-typed terms and values. The interpreter
of a closed program $M$ is $k \mapsto \mathit{eval}_k(\emptyset, M)$.

\paragraph{Interpreter for x86$\mu$}

A machine state $m$ of x86$\mu$ has the registers $\kw{rax}$,
$\kw{rdi}$, and $\kw{rsi}$, each holding a value or unset; the frame
$\sigma$, which holds a value in each slot assigned so far; the
environment of the open closure, if any; and the flag, which is a
Boolean or unset. An operand means $\mathit{arg}(a, m)$: the value of
an atom in the frame, the value of a register, and
$\kw{fold}\,(\mathit{arg}(a, m))$ for $\kw{fold}\,a$ and the $v$ such
that $\mathit{arg}(a, m) = \kw{fold}\,v$ for $\kw{unfold}\,a$. An
object means the value that its operands make:
$\mathit{obj}(\kw{pair}(a_1, a_2), m) = \kw{pack}(\mathit{arg}(a_1, m),
\mathit{arg}(a_2, m))$, and likewise a tuple and the injections. A write
$\mathit{put}(d, v, m)$ sets a register, or, for a slot $x$, extends
the frame cut at $x$ with $v$, $\sigma|_x[x{:=}v]$, and closes the open
closure, as in the graph-model semantics. The instructions update the
state.
\defhead{The effect of an instruction}{\mathit{instr}(i, m) = m'}
\begin{align*}
  \kw{movq}\; a, d &: \mathit{put}(d, \mathit{arg}(a, m), m) \\
  \kw{subq}\; a, \%r &: \mathit{put}(\%r, \delta_{\kw{sub}}(m(r), \mathit{arg}(a, m)), m) \\
  \kw{cmpq}\; a_2, a_1 &: m \text{ with the flag } (\mathit{arg}(a_1, m) = \mathit{arg}(a_2, m)) \\
  \kw{sete}\; d &: \mathit{put}(d, \kw{inl}\,\kw{tup}() \text{ if the flag is true, else } \kw{inr}\,\kw{tup}(), m) \\
  \kw{alloc}\; o, d &: \mathit{put}(d, \mathit{obj}(o, m), m) \\
  \kw{load}\; \kw{get}_i\; a, d &: \mathit{put}(d, v_i, m) \text{ where } \mathit{arg}(a, m) = \kw{tup}(v_1, \ldots, v_n) \\
  \kw{fst}\; x, \%r &: m \text{ with } r \text{ set to } c \text{ and the open closure } u,
    \text{ where } \sigma(x) = \kw{pack}(c, u) \\
  \kw{snd}\; x, \%r &: m \text{ with } r \text{ set to the environment of the open closure}
\end{align*}
where $\kw{fst}$ also unsets the registers that held the code or the
environment of an earlier open closure. Write $\mathit{instr}(i, m)$
for the state that an instruction $i$ other than $\kw{callq}$ leaves.
The interpreter $\mathit{run}_k(B, m, s)$ runs a sequence $s$ with
blocks $B$ in the state $m$.
\defhead{The interpreter of x86$\mu$}{\mathit{run}_k(B, m, s) = e \qquad \mathit{jump}_k(B, l, \sigma) = e \qquad \mathit{call}_k(c, u, w) = e}
\begin{align*}
  \mathit{run}_k(B, m, i;\; s) &= \mathit{run}_k(B, \mathit{instr}(i, m), s)
    \qquad\text{for $i$ not $\kw{callq}$} \\
  \mathit{run}_k(B, m, \kw{callq}\; a;\; s) &= v \leftarrow \mathit{call}_k(\mathit{arg}(a, m), m(\kw{rdi}), m(\kw{rsi}));\;
    \mathit{run}_k(B, m_v, s) \\
  \mathit{run}_k(B, m, \kw{retq}) &= \mathsf{ret}\; m(\kw{rax}) \\
  \mathit{run}_k(B, m, \kw{jmp}\;l) &= \mathit{jump}_k(B, l, \sigma) \\
  \mathit{run}_k(B, m, \kw{case}\;x\;(y.\,l_1)\;(y.\,l_2)) &= \begin{cases}
      \mathit{jump}_k(B, l_1, \sigma|_y[y{:=}v]) & \text{if } \sigma(x) = \kw{inl}\,v\\
      \mathit{jump}_k(B, l_2, \sigma|_y[y{:=}v]) & \text{if } \sigma(x) = \kw{inr}\,v
    \end{cases} \\[4pt]
  \mathit{jump}_k(B, l, \sigma) &= \mathit{run}_k(B_l, m_\sigma, s_l) \\
  \mathit{call}_0(\kw{code}\,f, u, w) &= \mathsf{timeout} \\
  \mathit{call}_{k+1}(\kw{code}\,f, u, w) &= \mathit{run}_k(B_f, m_{u,w}, s_f)
\end{align*}
Here $\sigma$ is the frame of $m$; $m_v$ has the frame of $m$, $v$ in
$\kw{rax}$, and nothing else set; $s_l$ is block $l$ and $B_l$ the
blocks older than it; $m_\sigma$ has the frame $\sigma$ and nothing else
set; $s_f$ is the start of the function body of definition $f$, with
blocks $B_f$; and $m_{u,w}$ has $u$ in $\kw{rdi}$, $w$ in $\kw{rsi}$, an
empty frame, and nothing else set. The interpreter of a program is $k
\mapsto \mathit{run}_k(B, m_0, s)$ for the main function body $s$ with
blocks $B$, where $m_0$ has an empty frame and nothing set.

\begin{theorem}[Correctness, clocked interpreters]\label{thm:interp}
Let $M$ be a closed term of \iswim\ of type $\kw{Int}$. For every integer $n$, the interpreter of \iswim\ on $M$ terminates with $n$ if and only if the interpreter of x86$\mu$ on $\pass{compile}(M)$ does, and one diverges if and only if the other does. More generally, for every closed term $M$, the two interpreters are equivalent.
\end{theorem}

\subsection{Abstract Machines}
\label{app:machines}
\agda{Mu/Semantics/AbstractMachines/ISWIM.agda, Mu/Semantics/AbstractMachines/X86.agda}

\paragraph{Machine for \iswim}

A state of the CEK machine of \iswim\ either evaluates a term $M$ in an environment $\rho$ with a continuation $\kappa$, $\mathit{eval}(M, \rho, \kappa)$, or returns a value $v$ to a continuation, $\mathit{ret}(v, \kappa)$. A continuation is a stack of frames, and each frame holds the work that remains after a subterm returns its value; for example, $\mathit{appL}(M, \rho, \kappa)$ holds the argument $M$ that the machine evaluates next.
\defhead{States and continuations}{\varsigma \qquad \kappa}
\begin{align*}
  \varsigma &::= \mathit{eval}(M, \rho, \kappa) \mid \mathit{ret}(v, \kappa) \\
  \kappa &::= \mathit{halt} \mid \mathit{appL}(M, \rho, \kappa) \mid \mathit{appR}(f, \kappa)
    \mid \mathit{tupK}(\vec{v}, \vec{M}, \rho, \kappa) \mid \mathit{getK}_i(\kappa)
    \mid \mathit{inlK}(\kappa) \mid \mathit{inrK}(\kappa) \\
 &\mid \mathit{caseK}(x.\,M, y.\,N, \rho, \kappa) \mid \mathit{primL}_p(M, \rho, \kappa)
    \mid \mathit{primR}_p(n, \kappa) \\
 &\mid \mathit{letK}(x.\,N, \rho, \kappa) \mid \mathit{foldK}(\kappa) \mid \mathit{unfoldK}(\kappa)
\end{align*}
\defhead{Steps of the CEK machine of \iswim}{\varsigma \longrightarrow \varsigma'}
\begin{align*}
  \mathit{eval}(x, \rho, \kappa) &\longrightarrow \mathit{ret}(\rho(x), \kappa) \\
  \mathit{eval}(\lambda x.\,N, \rho, \kappa) &\longrightarrow \mathit{ret}(\langle \lambda x.\,N, \rho\rangle, \kappa) \\
  \mathit{eval}(L\;M, \rho, \kappa) &\longrightarrow \mathit{eval}(L, \rho, \mathit{appL}(M, \rho, \kappa)) \\
  \mathit{eval}(c, \rho, \kappa) &\longrightarrow \mathit{ret}(c, \kappa) \\
  \mathit{eval}(\kw{tuple}\,(), \rho, \kappa) &\longrightarrow \mathit{ret}(\kw{tup}(), \kappa) \\
  \mathit{eval}(\kw{tuple}\,(M_1, \ldots, M_n), \rho, \kappa) &\longrightarrow \mathit{eval}(M_1, \rho, \mathit{tupK}((), (M_2, \ldots, M_n), \rho, \kappa)) \\
  \mathit{eval}(\kw{get}_i\,M, \rho, \kappa) &\longrightarrow \mathit{eval}(M, \rho, \mathit{getK}_i(\kappa))
    \qquad\text{(likewise $\kw{inl}$, $\kw{inr}$, $\kw{fold}$, $\kw{unfold}$)} \\
  \mathit{eval}(\kw{case}\;L\;(x.\,M)\;(y.\,N), \rho, \kappa) &\longrightarrow \mathit{eval}(L, \rho, \mathit{caseK}(x.\,M, y.\,N, \rho, \kappa)) \\
  \mathit{eval}(\kw{prim}_p\;L\;M, \rho, \kappa) &\longrightarrow \mathit{eval}(L, \rho, \mathit{primL}_p(M, \rho, \kappa)) \\
  \mathit{eval}(\kw{let}\;x = M\;\kw{in}\;N, \rho, \kappa) &\longrightarrow \mathit{eval}(M, \rho, \mathit{letK}(x.\,N, \rho, \kappa)) \\[4pt]
  \mathit{ret}(f, \mathit{appL}(M, \rho, \kappa)) &\longrightarrow \mathit{eval}(M, \rho, \mathit{appR}(f, \kappa)) \\
  \mathit{ret}(w, \mathit{appR}(\langle \lambda x.\,N, \rho\rangle, \kappa)) &\longrightarrow \mathit{eval}(N, \rho[x{:=}w], \kappa) \\
  \mathit{ret}(v, \mathit{tupK}(\vec{v}, (), \rho, \kappa)) &\longrightarrow \mathit{ret}(\kw{tup}(\vec{v}, v), \kappa) \\
  \mathit{ret}(v, \mathit{tupK}(\vec{v}, (M, \vec{M}), \rho, \kappa)) &\longrightarrow \mathit{eval}(M, \rho, \mathit{tupK}((\vec{v}, v), \vec{M}, \rho, \kappa)) \\
  \mathit{ret}(\kw{tup}(v_1, \ldots, v_n), \mathit{getK}_i(\kappa)) &\longrightarrow \mathit{ret}(v_i, \kappa) \\
  \mathit{ret}(v, \mathit{inlK}(\kappa)) &\longrightarrow \mathit{ret}(\kw{inl}\,v, \kappa)
    \qquad\text{(likewise $\kw{inr}$, $\kw{fold}$)} \\
  \mathit{ret}(\kw{inl}\,v, \mathit{caseK}(x.\,M, y.\,N, \rho, \kappa)) &\longrightarrow \mathit{eval}(M, \rho[x{:=}v], \kappa) \\
  \mathit{ret}(\kw{inr}\,v, \mathit{caseK}(x.\,M, y.\,N, \rho, \kappa)) &\longrightarrow \mathit{eval}(N, \rho[y{:=}v], \kappa) \\
  \mathit{ret}(n, \mathit{primL}_p(M, \rho, \kappa)) &\longrightarrow \mathit{eval}(M, \rho, \mathit{primR}_p(n, \kappa)) \\
  \mathit{ret}(m, \mathit{primR}_p(n, \kappa)) &\longrightarrow \mathit{ret}(\delta_p(n, m), \kappa) \\
  \mathit{ret}(v, \mathit{letK}(x.\,N, \rho, \kappa)) &\longrightarrow \mathit{eval}(N, \rho[x{:=}v], \kappa) \\
  \mathit{ret}(\kw{fold}\,v, \mathit{unfoldK}(\kappa)) &\longrightarrow \mathit{ret}(v, \kappa)
\end{align*}
\defhead{Final observation of the CEK machine of \iswim}{\varsigma \longrightarrow \mathsf{stop}(o)}
\[
  \mathit{ret}(v, \mathit{halt}) \longrightarrow \mathsf{stop}(\mathit{obs}(v))
\]
The machine of a closed program $M$ starts in
$\mathit{eval}(M, \emptyset, \mathit{halt})$.

\paragraph{Machine for x86$\mu$}

In the notation of Appendix~\ref{app:interp}, a state of the machine of x86$\mu$, $\mathit{run}(s, B, \Sigma, m)$, runs a sequence $s$ with blocks $B$, a stack $\Sigma$ of the callers' return points, and registers, frame, open closure and flag $m$. A return point $\mathit{retK}(s, B, \sigma)$ holds the rest $s$ of the caller's sequence, its blocks, and its frame, which the callee's $\kw{retq}$ restores.
\defhead{States and stacks}{\varsigma \qquad \Sigma}
\begin{align*}
  \varsigma &::= \mathit{run}(s, B, \Sigma, m) \\
  \Sigma &::= \mathit{halt} \mid \mathit{retK}(s, B, \sigma) :: \Sigma
\end{align*}
\defhead{Steps of the machine of x86$\mu$}{\varsigma \longrightarrow \varsigma'}
\begin{align*}
  \mathit{run}(i;\; s, B, \Sigma, m) &\longrightarrow \mathit{run}(s, B, \Sigma, \mathit{instr}(i, m))
    \qquad\text{for $i$ not $\kw{callq}$} \\
  \mathit{run}(\kw{callq}\; a;\; s, B, \Sigma, m) &\longrightarrow \mathit{run}(s_f, B_f, \mathit{retK}(s, B, \sigma) :: \Sigma, m_{u,w}) \\
 & \quad\text{where } \mathit{arg}(a, m) = \kw{code}\,f,\ u = m(\kw{rdi}),\ w = m(\kw{rsi}) \\
  \mathit{run}(\kw{retq}, B, \mathit{retK}(s, B', \sigma') :: \Sigma, m) &\longrightarrow \mathit{run}(s, B', \Sigma, m') \\
 & \quad\text{where $m'$ has the frame $\sigma'$ and only $\kw{rax}$ set, to $m(\kw{rax})$} \\
  \mathit{run}(\kw{jmp}\;l, B, \Sigma, m) &\longrightarrow \mathit{run}(s_l, B_l, \Sigma, m_\sigma) \\
  \mathit{run}(\kw{case}\;x\;(y.\,l_1)\;(y.\,l_2), B, \Sigma, m) &\longrightarrow \mathit{run}(s_{l_1}, B_{l_1}, \Sigma, m_{\sigma|_y[y{:=}v]}) \\
 & \quad\text{if } \sigma(x) = \kw{inl}\,v \\
 & \text{(and likewise for $\kw{inr}$ and $l_2$)}
\end{align*}
\defhead{Final observation of the machine of x86$\mu$}{\varsigma \longrightarrow \mathsf{stop}(o)}
\[
  \mathit{run}(\kw{retq}, B, \mathit{halt}, m) \longrightarrow \mathsf{stop}(\mathit{obs}(m(\kw{rax})))
\]
The machine of a program starts in $\mathit{run}(s, B, \mathit{halt},
m_0)$ for the main function body $s$ with blocks $B$.

\begin{theorem}[Correctness, abstract machines]\label{thm:machines}
Let $M$ be a closed term of \iswim\ of type $\kw{Int}$. For every integer $n$, the machine of \iswim\ on $M$ terminates with $n$ if and only if the machine of x86$\mu$ on $\pass{compile}(M)$ does, and one diverges if and only if the other does. More generally, for every closed term $M$, the two machines are equivalent.
\end{theorem}

\subsection{Reduction Semantics}
\label{app:reduction}
\agda{Mu/Semantics/ReductionSemantics/ISWIM.agda, Framework/Reduction.agda}

The reduction semantics of \iswim\ rewrites closed terms, by
substitution, call by value and left to right. The values are
\[
  V, W ::= \lambda x.\,N \mid c \mid \kw{tuple}\,(V_1, \ldots, V_n) \mid \kw{inl}\,V
    \mid \kw{inr}\,V \mid \kw{fold}\,V
\]
The function $\mathit{red}$ maps a closed term $M$ to a \emph{result},
ranged over by $r$, which says whether $M$ is a value ($\mathsf{val}$)
or reduces to a term $M'$ ($\mathsf{to}\,M'$):
\[
  r ::= \mathsf{val} \mid \mathsf{to}\,M'
\]
In each case below, the subterms are reduced left to right. A subterm
that reduces makes the whole term reduce in the same position, and once
the subterms are values the term takes its own step.
\defhead{One reduction of a closed term}{\mathit{red}(M) = r}
\begin{align*}
  \mathit{red}(V) &= \mathsf{val} \\
  \mathit{red}(L\;M) &= \mathsf{to}\,(L'\;M) &&\text{if } \mathit{red}(L) = \mathsf{to}\,L' \\
 &= \mathsf{to}\,(V\;M') &&\text{if } L = V,\ \mathit{red}(M) = \mathsf{to}\,M' \\
 &= \mathsf{to}\,(N[x{:=}W]) &&\text{if } L = \lambda x.\,N,\ M = W \\
  \mathit{red}(\kw{tuple}\,(\vec{V}, M, \vec{M})) &= \mathsf{to}\,(\kw{tuple}\,(\vec{V}, M', \vec{M})) &&\text{if } \mathit{red}(M) = \mathsf{to}\,M' \\
  \mathit{red}(\kw{get}_i\,M) &= \mathsf{to}\,(\kw{get}_i\,M') &&\text{if } \mathit{red}(M) = \mathsf{to}\,M' \\
 &= \mathsf{to}\,V_i &&\text{if } M = \kw{tuple}\,(V_1, \ldots, V_n) \\
  \mathit{red}(\kw{inl}\,M) &= \mathsf{to}\,(\kw{inl}\,M') &&\text{if } \mathit{red}(M) = \mathsf{to}\,M' \\
 & &&\text{(likewise $\kw{inr}$ and $\kw{fold}$)} \\
  \mathit{red}(\kw{case}\;L\;(x.\,M)\;(y.\,N)) &= \mathsf{to}\,(\kw{case}\;L'\;(x.\,M)\;(y.\,N)) &&\text{if } \mathit{red}(L) = \mathsf{to}\,L' \\
 &= \mathsf{to}\,(M[x{:=}V]) &&\text{if } L = \kw{inl}\,V \\
 &= \mathsf{to}\,(N[y{:=}V]) &&\text{if } L = \kw{inr}\,V \\
  \mathit{red}(\kw{prim}_p\;L\;M) &= \mathsf{to}\,(\kw{prim}_p\;L'\;M) &&\text{if } \mathit{red}(L) = \mathsf{to}\,L' \\
 &= \mathsf{to}\,(\kw{prim}_p\;V\;M') &&\text{if } L = V,\ \mathit{red}(M) = \mathsf{to}\,M' \\
 &= \mathsf{to}\,\delta_p(n, m) &&\text{if } L = n,\ M = m \\
  \mathit{red}(\kw{let}\;x{=}M\;\kw{in}\;N) &= \mathsf{to}\,(\kw{let}\;x{=}M'\;\kw{in}\;N) &&\text{if } \mathit{red}(M) = \mathsf{to}\,M' \\
 &= \mathsf{to}\,(N[x{:=}V]) &&\text{if } M = V \\
  \mathit{red}(\kw{unfold}\;M) &= \mathsf{to}\,(\kw{unfold}\;M') &&\text{if } \mathit{red}(M) = \mathsf{to}\,M' \\
 &= \mathsf{to}\,V &&\text{if } M = \kw{fold}\,V
\end{align*}
Here $\delta_p(n, m)$ is read as a term. The result of a primitive is a
value, and the result $\mathit{obs}(V)$ of a value is as in
Appendix~\ref{app:observations}, with $\mathit{fun}$ for $\lambda
x.\,N$. A closed well-typed term is a value or reduces, so these cases
cover every such term.\footnote{The Agda definition also has a result
$\mathsf{error}$ for a term that is stuck, which no closed well-typed
term is, and it writes the cases that reduce a subterm with one
combinator, $\mathit{step\text{-}in}$, that propagates a subterm's
reduction to the whole term. The equations above are that definition
with the combinator unfolded.} The reduction semantics of a closed
program $M$ is the machine whose states are closed terms and which
starts at $M$.
\defhead{Steps and final observation of the reduction semantics}{M \longrightarrow M' \qquad M \longrightarrow \mathsf{stop}(o)}
\[
  M \longrightarrow M' \ \text{ if } \mathit{red}(M) = \mathsf{to}\,M'
  \qquad\qquad
  M \longrightarrow \mathsf{stop}(\mathit{obs}(M)) \ \text{ if } \mathit{red}(M) = \mathsf{val}
\]

The first-order languages Imp, CFG, and x86$\mu$ have no binders, so
the reduction-semantics development uses the machines of
Appendix~\ref{app:machines} for them, and shares the proofs of
\pass{explicate control} and \pass{select instructions} with the
abstract-machine development.

\begin{theorem}[Correctness, reduction semantics]\label{thm:reduction}
Let $M$ be a closed term of \iswim\ of type $\kw{Int}$. For every integer $n$, the reduction semantics of \iswim\ on $M$ terminates with $n$ if and only if the machine of x86$\mu$ on $\pass{compile}(M)$ does, and one diverges if and only if the other does. More generally, for every closed term $M$, the two are equivalent.
\end{theorem}

\subsection{Scott Domains}
\label{app:scott}
\agda{Framework/Scott/Lift.agda, Framework/Scott/Domain.agda, Framework/Scott/Ops.agda, Mu/Semantics/ScottDomains/ISWIM.agda, Mu/Semantics/ScottDomains/X86.agda}

Our Agda development builds the domain of values with the constructive
domain theory of TypeTopology \citep{deJong:2023aa}, in which domains
are directed-complete partial orders, as the bilimit of
Appendix~\ref{app:scott-delay}. The domain of values is the solution of
\[
  \mathbb{D} \cong \mathcal{L}(\Sigma_{t \in \mathit{Tag}}\, \mathit{Payload}_t(\mathbb{D}))
\]
where $\mathcal{L}$ is lifting, the standard construction that adds a
least element $\bot$ below a domain \citep{Winskel:1993aa,Escardo:2017ab}. The tag $t$ of a value distinguishes constants, tuples of each arity,
injections, functions, and pairs, and the payload holds its contents:
\[
\begin{array}{@{}r@{\;}c@{\;}l@{}}
  t & ::= & c \mid \mathit{tup}_n \mid \mathit{left} \mid \mathit{right}
    \mid \mathit{fun} \mid \mathit{pr}
\end{array}
\]
\[
\begin{array}{@{}r@{\;}c@{\;}l@{\qquad}r@{\;}c@{\;}l@{}}
  \mathit{Payload}_c(\mathbb{D}) & = & \mathbf{1} &
  \mathit{Payload}_{\mathit{fun}}(\mathbb{D}) & = & [\mathbb{D} \to \mathbb{D}]\\
  \mathit{Payload}_{\mathit{tup}_n}(\mathbb{D}) & = & \mathbb{D}^n &
  \mathit{Payload}_{\mathit{pr}}(\mathbb{D}) & = & \mathbb{D}^2\\
  \mathit{Payload}_{\mathit{left}}(\mathbb{D}) & = & \mathbb{D} &
  \mathit{Payload}_{\mathit{right}}(\mathbb{D}) & = & \mathbb{D}
\end{array}
\]
where $\mathbf{1}$ is the one-point domain and $[\mathbb{D} \to \mathbb{D}]$
the domain of continuous functions. The least element $\bot$ of $\mathbb{D}$ means divergence;
a value $d$ is \emph{defined} when $d \neq \bot$, that is, when
$d = \up{(t, x)}$ for a tag $t$ and a payload $x$ in
$\mathit{Payload}_t(\mathbb{D})$, where $\up{\cdot}$ is the injection of
lifting. The empty tuple is $\langle\rangle = \up{(\mathit{tup}_0, ())}$,
where $()$ is the element of $\mathbf{1}$ and of $\mathbb{D}^0$. As in the graph model, the operators are
strict, and each of them is also continuous. We define them as follows,
where $g'$ is the strict function that agrees with $g$ on defined
values and maps $\bot$ to $\bot$.
\begin{align*}
  \Lambda(g) &= \up{(\mathit{fun}, g')} \\
  \mathsf{app}(d, e) &= h(e) &&\text{if } d = \up{(\mathit{fun}, h)} \text{ and } e \text{ is defined} \\
  \mathcal{T}(d_1, \ldots, d_n) &= \up{(\mathit{tup}_n, (d_1, \ldots, d_n))} &&\text{if every } d_i \text{ is defined} \\
  \pi_{\mathit{tup}_i}(d) &= d_i &&\text{if } d = \up{(\mathit{tup}_n, (d_1, \ldots, d_n))} \\
  \mathsf{pair}(d, e) &= \up{(\mathit{pr}, (d, e))} &&\text{if } d \text{ and } e \text{ are defined} \\
  \car(d) &= d_1 &&\text{if } d = \up{(\mathit{pr}, (d_1, d_2))} \\
  \cdr(d) &= d_2 &&\text{if } d = \up{(\mathit{pr}, (d_1, d_2))} \\
  \mathsf{inl}(d) &= \up{(\mathit{left}, d)} &&\text{if } d \text{ is defined} \\
  \mathsf{inr}(d) &= \up{(\mathit{right}, d)} &&\text{if } d \text{ is defined} \\
  \mathsf{cases}(d, f, g) &= f(e) &&\text{if } d = \up{(\mathit{left}, e)} \\
   &= g(e) &&\text{if } d = \up{(\mathit{right}, e)} \\
  \mathsf{prim}_p(d, e) &= \mathit{out}(\delta(p, n, m)) &&\text{if } d = \up{(n, ())} \text{ and } e = \up{(m, ())} \text{ for integers } n, m \\
  \mathsf{seq}(d, e) &= e &&\text{if } d \text{ is defined}
\end{align*}
In every case that the equations do not cover, the result is $\bot$.
As in the graph model (Appendix~\ref{app:operators}), $\mathit{out}$ turns
the result of a primitive into a value: $\mathit{out}(k) = \up{(k, ())}$
for an integer $k$, $\mathit{out}(\mathsf{true}) = \up{(\mathit{left}, \langle\rangle)}$,
and $\mathit{out}(\mathsf{false}) = \up{(\mathit{right}, \langle\rangle)}$.
An environment maps variables to values, and the environment $\rho_0$
of a closed program maps every variable to the empty tuple
$\langle\rangle$.

\paragraph{Semantics of \iswim}

The semantics of \iswim\ has the equations of the graph-model semantics
(Appendix~\ref{app:languages}), with the operators of $\mathbb{D}$:
\[
\begin{array}{@{}r@{\;}c@{\;}l@{\qquad}r@{\;}c@{\;}l@{}}
  \den{x}\rho & = & \rho(x) &
  \den{c}\rho & = & \up{(c, ())}\\
  \den{\lambda x.\, N}\rho & = & \Lambda(\mlam{d} \den{N}\rho[x{:=}d]) &
  \den{L\;M}\rho & = & \mathsf{app}(\den{L}\rho, \den{M}\rho)\\
  \den{\kw{tuple}\;(M_1, \ldots, M_n)}\rho & = & \mathcal{T}(\den{M_1}\rho, \ldots, \den{M_n}\rho) &
  \den{\kw{get}_i\,M}\rho & = & \pi_{\mathit{tup}_i}(\den{M}\rho)\\
  \den{\kw{inl}\;M}\rho & = & \mathsf{inl}(\den{M}\rho) &
  \den{\kw{inr}\;M}\rho & = & \mathsf{inr}(\den{M}\rho)\\
  \den{\kw{prim}_p\;L\;M}\rho & = & \mathsf{prim}_p(\den{L}\rho, \den{M}\rho) &
  \den{\kw{fold}\;M}\rho & = & \den{M}\rho\\
  \den{\kw{let}\;x = M\;\kw{in}\;N}\rho & = & \mathsf{app}(\Lambda(\mlam{d} \den{N}\rho[x{:=}d]), \den{M}\rho) &
  \den{\kw{unfold}\;M}\rho & = & \den{M}\rho
\end{array}
\]
\[
  \den{\kw{case}\;L\;(x.\,M)\;(y.\,N)}\rho =
    \mathsf{cases}(\den{L}\rho,\; (\mlam{d} \den{M}\rho[x{:=}d])',\; (\mlam{d} \den{N}\rho[y{:=}d])')
\]
where $g'$ is the strict function that agrees with $g$ on defined
values, as in $\Lambda$.

\paragraph{Semantics of x86$\mu$}

The semantics of x86$\mu$ has the equations of the graph-model
semantics of x86$\mu$ (Appendix~\ref{app:languages}), with the
operators of $\mathbb{D}$ in place of those of the graph model,
$\mathsf{seq}$ in place of $\mathsf{guard}$ in the meaning of
$\kw{callq}$, and the empty tuple $\langle\rangle$ in place of
$\{\langle\rangle\}$ as the content of an unset register and an unset
flag. A sequence means a continuous function of the frame, the
registers, and the flag.

\begin{theorem}[Correctness, Scott domains]\label{thm:scott} Assume propositional truncation, function extensionality, and propositional extensionality. Let $M$ be a closed term of \iswim\ of type $\kw{Int}$. For every constant $c$, $\den{M}\rho_0 = \up{(c, ())}$ if and only if $\den{\pass{compile}(M)} = \up{(c, ())}$.
Moreover, $\den{M}\rho_0$ is defined if and only if
$\den{\pass{compile}(M)}$ is defined. More generally, for every closed
term $M$, the same holds.
\end{theorem}

\section{The Scott-Domain Proof of \pass{delay}}
\label{app:scott-delay}
\agda{Mu/proof/Correctness/ScottDomains/DelayRel.agda, DelayPreserve.agda, DelayReflect.agda, DelayTerms.agda, DelayCorrect.agda; proof/Framework/Scott/Stages.agda, StageCall.agda}

As in the graph model (Section~\ref{sec:approx-lr}), \pass{delay}
changes the representation of closures, so the Scott-domain proof
relates the meaning of a program to the meaning of its translation by
a logical relation, one for each direction. The universal domain
$\mathbb{D}$ (Appendix~\ref{app:scott}) rules out the two usual ways
of defining such a relation. Recursion on types does not terminate,
because the types include recursive types, and recursion on values is
not well founded, because values are infinite. We instead define the
relations by recursion on the stages of the bilimit that constructs
$\mathbb{D}$, as \citet{Reynolds:1974aa} did. Each relation relates an
approximation of one value at a stage to the whole other value, so,
like the approximant logical relation, it approximates only one side.
The preservation relation approximates the source, and the reflection
relation approximates the target.

\paragraph{Stages}

The stages of $\mathbb{D}$ are the domains
\[
  \mathbb{D}_0 = \mathbf{1}
  \qquad
  \mathbb{D}_{n+1} = \mathcal{L}(\Sigma_{t \in \mathit{Tag}}\, \mathit{Payload}_t(\mathbb{D}_n))
\]
with the maps $\pi_n : \mathbb{D}_{n+1} \to \mathbb{D}_n$ and
$\varepsilon_n : \mathbb{D}_n \to \mathbb{D}_{n+1}$ between consecutive
stages, and the maps $\pi^\infty_n : \mathbb{D} \to \mathbb{D}_n$ and
$\varepsilon^\infty_n : \mathbb{D}_n \to \mathbb{D}$ between a stage and
the limit. Every value $d$ is the supremum of its approximations
$\varepsilon^\infty_n(\pi^\infty_n(d))$. At stage $n + 1$, the payload of a
function is in $[\mathbb{D}_n \to \mathbb{D}_n]$ and the payload of a
pair is in $\mathbb{D}_n^2$.

A target closure is a pair $(c, e)$ of code and environment, and the
meaning of $\kw{callc}$ calls it on an argument $a$ by
$\mathit{call}(c, e, a) = \mathsf{app}(\mathsf{app}(c, e), a)$. The
code is a curried function that the call applies twice, and applying a
function of stage $k + 1$ gives a result of stage $k$. A call of a
pair of stage $m + 2$ therefore gives its result at stage $m$:
\[
  \mathit{call}_m((c, e), v) =
    \pi^\infty_m(\mathit{call}(\varepsilon^\infty_{m+2}(c),\;
      \varepsilon^\infty_{m+1}(\pi_{m+1}(e)),\;
      \varepsilon^\infty_m(v)))
  \qquad \text{for } c, e \in \mathbb{D}_{m+2} \text{ and } v \in \mathbb{D}_m.
\]

\begin{figure}[h]
\small\figdefs
\defhead{Related tags}{t \SameTag{+} t' \qquad t \SameTag{-} t'}
\[
  t \SameTag{\pm} t \text{ for } t \notin \{\mathit{fun}, \mathit{pr}\}
  \qquad
  \mathit{fun} \SameTag{+} \mathit{pr}
  \qquad
  \mathit{pr} \SameTag{-} \mathit{fun}
\]
\defhead{The stage relations}{s \SR{\pm}{n} w \quad (s \in \mathbb{D}_n,\ w \in \mathbb{D})}
\[
\begin{array}{@{}r@{\;}c@{\;}l@{}}
  s \SR{\pm}{0} w & \iff & \mathsf{true}\\
  s \SR{\pm}{n+1} w & \iff &
    \forall t, x.\; s = \up{(t, x)} \Rightarrow
    \exists t', y.\; w = \up{(t', y)} \land t \SameTag{\pm} t'
    \land (t, x) \mathrel{\mathcal{M}^{\pm}_n} (t', y)
\end{array}
\]
\defhead{Related payloads}{(t, x) \mathrel{\mathcal{M}^{\pm}_n} (t', y) \quad (x \in \mathit{Payload}_t(\mathbb{D}_n),\ y \in \mathit{Payload}_{t'}(\mathbb{D}))}
\[
\begin{array}{@{}r@{\;}c@{\;}l@{}}
  (c, ()) \mathrel{\mathcal{M}^{\pm}_n} (c, ()) & \iff & \mathsf{true}\\
  (\mathit{tup}_k, (x_1, \ldots, x_k)) \mathrel{\mathcal{M}^{\pm}_n} (\mathit{tup}_k, (y_1, \ldots, y_k))
    & \iff & \forall i.\; x_i \SR{\pm}{n} y_i\\
  (\mathit{left}, x) \mathrel{\mathcal{M}^{\pm}_n} (\mathit{left}, y)
    & \iff & x \SR{\pm}{n} y
    \qquad \text{(likewise $\mathit{right}$)}\\
  (\mathit{fun}, h) \mathrel{\mathcal{M}^{+}_n} (\mathit{pr}, (c, e))
    & \iff & \forall v \in \mathbb{D}_n,\, a \in \mathbb{D}.\;
      v \SR{+}{n} a \Rightarrow h(v) \SR{+}{n} \mathit{call}(c, e, a)\\
  (\mathit{pr}, (c, e)) \mathrel{\mathcal{M}^{-}_{m+2}} (\mathit{fun}, g)
    & \iff & \forall v \in \mathbb{D}_m,\, a \in \mathbb{D}.\;
      v \SR{-}{m} a \Rightarrow \mathit{call}_m((c, e), v) \SR{-}{m} g'(a)\\
  (\mathit{pr}, (c, e)) \mathrel{\mathcal{M}^{-}_n} (\mathit{fun}, g)
    & \iff & \mathsf{true} \qquad \text{for } n < 2
\end{array}
\]
\defhead{Related values}{d \SRel{\pm} w \qquad \mathit{Pres}(d, d') \qquad \mathit{Refl}(d, d')}
\[
  d \SRel{\pm} w \iff \forall n.\; \pi^\infty_n(d) \SR{\pm}{n} w
  \qquad
  \mathit{Pres}(d, d') \iff d \SRel{+} d'
  \qquad
  \mathit{Refl}(d, d') \iff d' \SRel{-} d
\]
\caption{The relations of the Scott-domain proof of \pass{delay},
  between a source (EnvClos) value $d$ and a target (FlatClos) value
  $d'$. The superscript $\pm$ stands for either direction. Here $g'$ is
  the strict function that agrees with $g$ on defined values.}
\label{fig:scott-delay}
\end{figure}

\paragraph{The relations}

Figure~\ref{fig:scott-delay} defines the relations. At stage $n + 1$,
a defined approximated value $s$ is related to $w$ when $w$ is defined
too, the two have related tags, and their payloads are related at
stage $n$; an undefined approximated value is related to every value.
Related values have the same tag, except for closures. Under
preservation, a source function is related to a target pair when the
function and calls of the pair, on related arguments, give related
results; this clause has the shape of the relation of
\citet{Nielsen:2000aa} for defunctionalization. Under reflection, the
approximated value is the target pair, and its calls show their
results two stages down. The clause for a pair at stage $m + 2$
therefore speaks of stage $m$, and it says nothing at stages $0$ and
$1$. Two values are related when every approximation of the
approximated one is related to the other.

\paragraph{Properties of the stages}

The proof needs four properties of each stage relation, which hold by
induction on the stage.

\begin{lemma}[Stages]\label{lem:scott-stages}
For both directions and every $n$:
\begin{enumerate}
\item $s \SR{\pm}{n} w$ is a proposition;
\item it is downward closed in $s$, that is, if $s' \sqsubseteq s$ and
  $s \SR{\pm}{n} w$, then $s' \SR{\pm}{n} w$;
\item it is admissible in $s$, that is, if $s_i \SR{\pm}{n} w$ for every $i$
  of a directed family, then $\bigsqcup_i s_i \SR{\pm}{n} w$;
\item it survives moving between stages, that is, if $s \SR{\pm}{n+1} w$, then
  $\pi_n(s) \SR{\pm}{n} w$, and if $s \SR{\pm}{n} w$, then
  $\varepsilon_n(s) \SR{\pm}{n+1} w$.
\end{enumerate}
Consequently, if $s \SR{\pm}{n} w$, then
$\varepsilon^\infty_n(s) \SRel{\pm} w$.
\end{lemma}

The closure clauses make parts~3 and~4 the hard ones. For part~3 under
reflection, $\mathit{call}_m$ is continuous in the pair, so the call of
a supremum of pairs is the supremum of their calls. For part~4,
$\mathit{call}_m$ commutes with $\pi$ and $\varepsilon$ up to the order
of $\mathbb{D}_m$, and part~2 absorbs the difference.

\paragraph{The fundamental lemma}

Write $\sim$ for either $\mathit{Pres}$ or $\mathit{Refl}$. Both relate
$\bot$ to every value in the approximated position, and the operators
that EnvClos and FlatClos share preserve both.

\begin{lemma}[Operators]\label{lem:scott-ops}
If $d_i \sim d'_i$ for each $i$, then $\mathcal{T}(d_1, \ldots, d_n)
\sim \mathcal{T}(d'_1, \ldots, d'_n)$, and likewise for
$\pi_{\mathit{tup}_i}$, $\mathsf{inl}$, $\mathsf{inr}$,
$\mathsf{prim}_p$, and $\mathsf{seq}$; constants and $\langle\rangle$
are related to themselves. If the branches of $\mathsf{cases}$ map
related values to related values, then $\mathsf{cases}$ preserves
$\sim$.
\end{lemma}

The two languages differ only in closures and calls.

\begin{lemma}[Closures and calls]\label{lem:scott-clos}
Let $t \sim t'$, and let the bodies $B$ and $B'$ map related
environments and arguments to related values. Then
\[
  \mathsf{app}(\Lambda(\mlam{X} \Lambda(\mlam{Y} B(X, Y))), t)
  \sim
  \mathsf{pair}(\Lambda(\mlam{X} \Lambda(\mlam{Y} B'(X, Y))), t').
\]
If $l \sim l'$ and $a \sim a'$, then
$\mathsf{app}(l, a) \sim \mathsf{app}(\mathsf{app}(\car(l'), \cdr(l')), a')$.
\end{lemma}

Under preservation, both parts follow from the clause for functions
and part~4 of Lemma~\ref{lem:scott-stages}. Under reflection, the
closure case uses the clause for pairs. A call of the approximation at
stage $m + 2$ of the target pair is below the approximation at stage
$m$ of the call of the pair itself, and part~2 of
Lemma~\ref{lem:scott-stages} applies. In the call case, the
approximation at stage $n$ of the result of a call is the supremum of
the calls of the approximations of the code, the environment and the
argument, so part~3 reduces the call to the calls of approximations,
and each of those is below an approximation of a call at a larger
stage, of which the clause for pairs speaks.

\begin{lemma}[Fundamental lemma]\label{lem:scott-fundamental}
For each direction, let $M$ be a term of EnvClos and let $\rho$ and
$\rho'$ be environments related at every variable of $M$. Then
$\den{M}\rho \sim \den{\pass{delay}(M)}\rho'$.
\end{lemma}

The proof is by induction on $M$, with the same induction for both
directions. Lemma~\ref{lem:scott-ops} covers the shared operators, and
Lemma~\ref{lem:scott-clos} covers the closures and calls.

\begin{theorem}[Correctness of \pass{delay}, Scott domains]\label{thm:scott-delay}
Let $M$ be a closed term of EnvClos. For every constant $c$,
$\den{M}\rho_0 = \up{(c, ())}$ if and only if
$\den{\pass{delay}(M)}\rho_0 = \up{(c, ())}$. Moreover,
$\den{M}\rho_0$ is defined if and only if
$\den{\pass{delay}(M)}\rho_0$ is defined.
\end{theorem}

The empty environments are related, so Lemma~\ref{lem:scott-fundamental}
gives $\mathit{Pres}(\den{M}\rho_0, \den{\pass{delay}(M)}\rho_0)$ and
$\mathit{Refl}(\den{M}\rho_0, \den{\pass{delay}(M)}\rho_0)$. In either
relation, a defined approximated value has a defined partner, and a
defined value related to a constant is that constant. Preservation
gives the forward implications and reflection the backward ones.

